% =====================================================================
% Method.
% Adapted from RESEARCH_PLAN.md §2 (mask diagram), §7 (cost model),
% §4 (Δ non-inferiority margin), §6 axis P (position layouts).
% =====================================================================
% NOTE: This file is NOT included in preprint.tex. The active paper body is sections/main_story.tex.
% This file is retained for historical reference only.
\section{Method}
\label{sec:method}

\subsection{Mask family}
\label{sec:method-masks}

Let a prompt consist of $N$ documents of lengths $\ell_1, \dots, \ell_N$,
a query block of length $Q$, and an answer block. Write
$T = \sum_i \ell_i + Q + \text{answer length}$. Under standard causal
attention every token attends to every earlier token. Decoupled Document
Inference restricts this pattern to a block-diagonal structure over the
document region, parameterized by which query rows are allowed to bridge
the blocks.

Table~\ref{tab:mask-taxonomy} enumerates the mask variants in scope:
\texttt{full} (baseline), \texttt{ddi\_qf} (documents isolated, whole
query bridges), \texttt{ddi\_strict} (documents isolated, only the last
query row bridges), and \texttt{blind} (no query or answer access to documents,
retained as a control that verifies each metric measures retrieval).

\begin{table}[t]
\centering
\caption{Mask taxonomy (T1). Regenerated from the canonical
    \texttt{paper/evidence.json} manifest on every build.}
\label{tab:mask-taxonomy}
\IfFileExists{generated/T1_mask_taxonomy.tex}{%
    % Generated by paper/render.py from paper/evidence.json.
% Sources: ddi/masks.py sha256=2ee804f18a94f4bd709b7629d4962a7c47f732157284c4218f31239956ba7f54.
\begin{tabular}{lcccc}
\toprule
Variant & Doc$\leftrightarrow$Doc & Query$\leftrightarrow$Doc & Answer$\leftrightarrow$Doc & Causal \\
\midrule
\texttt{full} & yes & yes & yes & yes \\
\texttt{ddi\_qf} & no & yes (all rows) & yes & within block \\
\texttt{ddi\_strict} & no & yes (last row) & yes & within block \\
\texttt{ddi\_bridge\_k} & no & yes ($k$ rows) & yes & within block \\
\texttt{blind} & no & no & no & within block \\
\bottomrule
\end{tabular}
%
}{%
    \todoreport{T1 mask taxonomy missing}{./paper/build.sh}%
}
\end{table}

Figure~\ref{fig:mask-geometry} shows the information-flow change directly.
The crossed document edge is the removed interaction; query and answer access
remain available under \texttt{ddi\_qf}. This distinction prevents interpreting
DDI as merely dropping the retrieved context.

\begin{figure}[t]
\centering
\resizebox{\linewidth}{!}{% Generated by paper/render.py; aggregates: paper/evidence.json.
% Sources: VIBEQUEUE-53-REPORT.md sha256=7cd9223e128c5d82b61e59575035fe4bb54c9f48c5da721723fd5d214d85195e; results/cache_v53_local_20260907/summary.json sha256=28ae886d9dae851ed090b7ebb0744b33926060d9499c3f16a01221810a0e2189; results/cache_v53_local_20260907/timing.jsonl sha256=055cf47312269ca147c6d4517d26b82aabc910b11fd6ca07f5eba5e3f54153a8; results/cache_v53_local_20260907/validation.jsonl sha256=8597c0b5394b4b17248065ad1ebc386a0a2d864f089f8e65d51d5e293f7b8d4a; ddi/masks.py sha256=2ee804f18a94f4bd709b7629d4962a7c47f732157284c4218f31239956ba7f54
\begin{tikzpicture}[x=1cm,y=1cm,>=stealth,font=\small,box/.style={draw,very thick,rounded corners=3pt,text width=3.4cm,minimum height=1cm,align=center,inner sep=5pt}]
\node[anchor=west,text width=3.5cm,font=\bfseries,text=ddiBlue] at (0,4.3) {DDI-QF attention};
\fill[ddiBlue] (0.00,3.60) rectangle ++(.27,.27);
\fill[gray!9] (0.29,3.60) rectangle ++(.27,.27);
\fill[gray!9] (0.58,3.60) rectangle ++(.27,.27);
\fill[gray!9] (0.87,3.60) rectangle ++(.27,.27);
\fill[gray!9] (1.16,3.60) rectangle ++(.27,.27);
\fill[gray!9] (1.45,3.60) rectangle ++(.27,.27);
\fill[gray!9] (1.74,3.60) rectangle ++(.27,.27);
\fill[gray!9] (2.03,3.60) rectangle ++(.27,.27);
\fill[gray!9] (2.32,3.60) rectangle ++(.27,.27);
\fill[gray!9] (2.61,3.60) rectangle ++(.27,.27);
\fill[ddiBlue] (0.00,3.31) rectangle ++(.27,.27);
\fill[ddiBlue] (0.29,3.31) rectangle ++(.27,.27);
\fill[gray!9] (0.58,3.31) rectangle ++(.27,.27);
\fill[gray!9] (0.87,3.31) rectangle ++(.27,.27);
\fill[gray!9] (1.16,3.31) rectangle ++(.27,.27);
\fill[gray!9] (1.45,3.31) rectangle ++(.27,.27);
\fill[gray!9] (1.74,3.31) rectangle ++(.27,.27);
\fill[gray!9] (2.03,3.31) rectangle ++(.27,.27);
\fill[gray!9] (2.32,3.31) rectangle ++(.27,.27);
\fill[gray!9] (2.61,3.31) rectangle ++(.27,.27);
\fill[gray!9] (0.00,3.02) rectangle ++(.27,.27);
\fill[gray!9] (0.29,3.02) rectangle ++(.27,.27);
\fill[ddiBlue] (0.58,3.02) rectangle ++(.27,.27);
\fill[gray!9] (0.87,3.02) rectangle ++(.27,.27);
\fill[gray!9] (1.16,3.02) rectangle ++(.27,.27);
\fill[gray!9] (1.45,3.02) rectangle ++(.27,.27);
\fill[gray!9] (1.74,3.02) rectangle ++(.27,.27);
\fill[gray!9] (2.03,3.02) rectangle ++(.27,.27);
\fill[gray!9] (2.32,3.02) rectangle ++(.27,.27);
\fill[gray!9] (2.61,3.02) rectangle ++(.27,.27);
\fill[gray!9] (0.00,2.73) rectangle ++(.27,.27);
\fill[gray!9] (0.29,2.73) rectangle ++(.27,.27);
\fill[ddiBlue] (0.58,2.73) rectangle ++(.27,.27);
\fill[ddiBlue] (0.87,2.73) rectangle ++(.27,.27);
\fill[gray!9] (1.16,2.73) rectangle ++(.27,.27);
\fill[gray!9] (1.45,2.73) rectangle ++(.27,.27);
\fill[gray!9] (1.74,2.73) rectangle ++(.27,.27);
\fill[gray!9] (2.03,2.73) rectangle ++(.27,.27);
\fill[gray!9] (2.32,2.73) rectangle ++(.27,.27);
\fill[gray!9] (2.61,2.73) rectangle ++(.27,.27);
\fill[gray!9] (0.00,2.44) rectangle ++(.27,.27);
\fill[gray!9] (0.29,2.44) rectangle ++(.27,.27);
\fill[gray!9] (0.58,2.44) rectangle ++(.27,.27);
\fill[gray!9] (0.87,2.44) rectangle ++(.27,.27);
\fill[ddiBlue] (1.16,2.44) rectangle ++(.27,.27);
\fill[gray!9] (1.45,2.44) rectangle ++(.27,.27);
\fill[gray!9] (1.74,2.44) rectangle ++(.27,.27);
\fill[gray!9] (2.03,2.44) rectangle ++(.27,.27);
\fill[gray!9] (2.32,2.44) rectangle ++(.27,.27);
\fill[gray!9] (2.61,2.44) rectangle ++(.27,.27);
\fill[gray!9] (0.00,2.15) rectangle ++(.27,.27);
\fill[gray!9] (0.29,2.15) rectangle ++(.27,.27);
\fill[gray!9] (0.58,2.15) rectangle ++(.27,.27);
\fill[gray!9] (0.87,2.15) rectangle ++(.27,.27);
\fill[ddiBlue] (1.16,2.15) rectangle ++(.27,.27);
\fill[ddiBlue] (1.45,2.15) rectangle ++(.27,.27);
\fill[gray!9] (1.74,2.15) rectangle ++(.27,.27);
\fill[gray!9] (2.03,2.15) rectangle ++(.27,.27);
\fill[gray!9] (2.32,2.15) rectangle ++(.27,.27);
\fill[gray!9] (2.61,2.15) rectangle ++(.27,.27);
\fill[ddiGold] (0.00,1.86) rectangle ++(.27,.27);
\fill[ddiGold] (0.29,1.86) rectangle ++(.27,.27);
\fill[ddiGold] (0.58,1.86) rectangle ++(.27,.27);
\fill[ddiGold] (0.87,1.86) rectangle ++(.27,.27);
\fill[ddiGold] (1.16,1.86) rectangle ++(.27,.27);
\fill[ddiGold] (1.45,1.86) rectangle ++(.27,.27);
\fill[ddiGold] (1.74,1.86) rectangle ++(.27,.27);
\fill[gray!9] (2.03,1.86) rectangle ++(.27,.27);
\fill[gray!9] (2.32,1.86) rectangle ++(.27,.27);
\fill[gray!9] (2.61,1.86) rectangle ++(.27,.27);
\fill[ddiGold] (0.00,1.57) rectangle ++(.27,.27);
\fill[ddiGold] (0.29,1.57) rectangle ++(.27,.27);
\fill[ddiGold] (0.58,1.57) rectangle ++(.27,.27);
\fill[ddiGold] (0.87,1.57) rectangle ++(.27,.27);
\fill[ddiGold] (1.16,1.57) rectangle ++(.27,.27);
\fill[ddiGold] (1.45,1.57) rectangle ++(.27,.27);
\fill[ddiGold] (1.74,1.57) rectangle ++(.27,.27);
\fill[ddiGold] (2.03,1.57) rectangle ++(.27,.27);
\fill[gray!9] (2.32,1.57) rectangle ++(.27,.27);
\fill[gray!9] (2.61,1.57) rectangle ++(.27,.27);
\fill[ddiPurple] (0.00,1.28) rectangle ++(.27,.27);
\fill[ddiPurple] (0.29,1.28) rectangle ++(.27,.27);
\fill[ddiPurple] (0.58,1.28) rectangle ++(.27,.27);
\fill[ddiPurple] (0.87,1.28) rectangle ++(.27,.27);
\fill[ddiPurple] (1.16,1.28) rectangle ++(.27,.27);
\fill[ddiPurple] (1.45,1.28) rectangle ++(.27,.27);
\fill[ddiPurple] (1.74,1.28) rectangle ++(.27,.27);
\fill[ddiPurple] (2.03,1.28) rectangle ++(.27,.27);
\fill[ddiPurple] (2.32,1.28) rectangle ++(.27,.27);
\fill[gray!9] (2.61,1.28) rectangle ++(.27,.27);
\fill[ddiPurple] (0.00,0.99) rectangle ++(.27,.27);
\fill[ddiPurple] (0.29,0.99) rectangle ++(.27,.27);
\fill[ddiPurple] (0.58,0.99) rectangle ++(.27,.27);
\fill[ddiPurple] (0.87,0.99) rectangle ++(.27,.27);
\fill[ddiPurple] (1.16,0.99) rectangle ++(.27,.27);
\fill[ddiPurple] (1.45,0.99) rectangle ++(.27,.27);
\fill[ddiPurple] (1.74,0.99) rectangle ++(.27,.27);
\fill[ddiPurple] (2.03,0.99) rectangle ++(.27,.27);
\fill[ddiPurple] (2.32,0.99) rectangle ++(.27,.27);
\fill[ddiPurple] (2.61,0.99) rectangle ++(.27,.27);
\node[anchor=north west,text width=3.8cm,align=left] at (0,.8) {Keys $\rightarrow$; rows $\downarrow$\par Blue: documents\par Orange: query\par Purple: answer};
\node[anchor=west,text width=8.6cm,font=\bfseries,text=ddiGold] at (4.1,4.3) {COLD MISS: target vs implementation};
\node[box,draw=ddiGold,fill=ddiGold!9] (a) at (6,3.3) {Target: sparse prefill\par Parity gate pending};
\node[box,draw=ddiBlue,fill=ddiBlue!8] (b) at (10.5,3.3) {V53: dense SDPA\par per-document prefill};
\node[box,draw=ddiGreen,fill=ddiGreen!10] (c) at (6,1.6) {Persist KVs\par exact identity key};
\node[box,draw=ddiPurple,fill=ddiPurple!8] (d) at (10.5,1.6) {WARM HIT: load KVs\par Skip document prefill};
\draw[->,thick,dashed,ddiGold] (a) -- (c);
\draw[->,thick,ddiBlue] (b.south) -- ++(0,-.25) -| (c.north east);
\draw[->,thick,ddiGreen] (c) -- (d);
\node[anchor=west,text width=8.3cm,align=left] at (4.2,.05) {Dashed: unvalidated target kernel. Solid: implemented cache path. V53 does not measure a sparse kernel.};
\end{tikzpicture}
}
\caption{Information-flow intuition for full causal attention and DDI-QF.
    Arrows summarize permitted earlier-to-later information flow rather than
    individual token-level edges. DDI removes cross-document attention while
    preserving document integration through query and answer rows.}
\label{fig:mask-geometry}
\end{figure}

\subsection{Cost model}
\label{sec:method-cost}

For prefill, count lower-triangular attended (query,\,key) pairs and treat all
documents as equal length $\ell$. With $T=N\ell+Q$, $N>1$, and $Q>0$:
\begin{align}
P_{\mathrm{full}} &= \tfrac{T(T+1)}{2}, \label{eq:cost-full} \\
P_{\text{\texttt{ddi\_qf}}} &= \tfrac{N\ell(\ell+1)}{2}+QN\ell+\tfrac{Q(Q+1)}{2}, \label{eq:cost-qf} \\
P_{\text{\texttt{ddi\_strict}}} &= \tfrac{N\ell(\ell+1)}{2}+N\ell+\tfrac{Q(Q+1)}{2}. \label{eq:cost-strict}
\end{align}
The \texttt{ddi\_strict} expression includes its one last-query-row bridge;
answer rows are outside this prefill count and directly attend documents for
every non-blind variant. At $N=1$, the implementation deliberately returns the
full causal mask, so these isolation-family expressions are not a separate
semantic behavior. As $Q \to 0$, the full-to-block-diagonal pair-count ratio
approaches $N$. A larger query adds the $QN\ell$ bridge term and lowers that
ratio; pair counts are not wall-clock speedup predictions.

Section~\ref{sec:results-efficiency} compares these pair counts with retained
SDPA TTFT observations only at $\ell=512$, $Q=64$, and the listed $N$ values.
It does not identify a backend mechanism or report an $N$--$\ell$ timing grid.

\subsection{Position layouts and caching}
\label{sec:method-positions}

We consider two layouts for the RoPE position indices assigned to document
tokens:
\begin{itemize}
    \item \textbf{Sequential.} Documents occupy consecutive absolute
          positions, as in standard concatenation. A document's KV block is
          then a function of both its content and its position in the
          retrieved set; changing either invalidates the cache.
    \item \textbf{Restart.} Each document's positions restart at zero.
           Under a block-diagonal mask, a document's KV block can be a
           function of that document alone under the stated model and cache
           assumptions. This can make blocks keyable by document content hash,
           order-independent, and reusable across requests.
\end{itemize}
The restart layout is designed to turn DDI's structural sparsity into a
caching mechanism: under a block-diagonal mask, a document's KV block
can be a function of that document alone, making content-hash cache keys a
testable design consequence rather than a universal serving result.
The generated answer positions continue the query tail in either layout.
This is required for a cached decode to preserve the position scheme used for
prefill; Section~\ref{sec:results-caching-prereq} reports the corrected
restart control without a parity conclusion.

\subsection{Non-inferiority margin \texorpdfstring{$\Delta$}{Delta}}
\label{sec:method-delta}

Legacy tables do not provide fixed held-out IDs, paired outcomes, or independent
checkpoint replications, so they remain descriptive. Separately, the matched
three-seed transfer study retains fixed manifests, paired outcomes, and its
preregistered gates; its normalized remaining tuned gap is not a training
recovery estimand (Section~\ref{sec:results-quality}).

Concretely, $\Delta$ is computed as $k\cdot \sigma_{\text{full}}$ for a
fixed $k$ committed alongside the baseline it was derived from, using the
in-repo \texttt{ddi-delta} tool. Table~\ref{tab:t6-delta} reports the
resulting legacy baseline summaries. They are not comparison confidence
intervals or a parity band for the other legacy rows.
